\chapter{Pi, or ``Mmm... pie.''}

In an episode that first aired on March 11, 2001, Professor Frink finds 
himself unable to quiet down a rowdy audience of scientists. To
get their attention, he finally shouts out, ``Pi is exactly three!''

We now proceed to disprove this outrageous statement.\\

\begin{dfn}
A real number $r$ is \emph{transcendental} if there does not exist
any polynomial $P(x)$ with integer coefficients such that $P(r)=0$.\\
\end{dfn}

\begin{thm}\label{transcendentalpi} 
% \label allows us to refer back to this by number later by
%   using the \ref command.
The number $\pi$ is transcendental.\footnote{Notice how we avoided
the terrible sin of beginning a sentence with a number or symbol.}
\end{thm}

\begin{proof}
The proof of this theorem is just slightly beyond the scope of this 
thesis. It was first accomplished by Ferdinand von Lindemann 
in 1882. Interested readers should see his paper in \cite{FL1882}.\\
% \cite allows us to refer to items in the bibliography.
\end{proof}

\begin{thm}\label{3nottrans}
The number 3 is not transcendental.
\end{thm}

\begin{proof}
According to our definition above, it is sufficent to find a polynomial
$P(x)$ with all integer coefficients such that $P(3)=0$. 

Consider $P(x)=x^{132}-3x^{131}-2x^{74}+6x^{73}-x^2+9$. Clearly, $P(x)$ 
is a polynomial and its coefficients are all integers. And even a 
simpleton can see at a glance that $P(3)=0$.\\
\end{proof}

\begin{cor}
The number $\pi$ is not exactly 3.
\end{cor}

\begin{proof}
We saw in Theorem \ref{transcendentalpi} that $\pi$ is transcendental
and in Theorem \ref{3nottrans} that 3 is not transcendental. A number
can not be both transcendental and not transcendental\footnote{This
assumes that you believe in the Law of the Excluded Middle. If not, 
you may be in for a difficult time ahead.}, so $\pi$ must not be equal
 to 3.\\
\end{proof}

We shouldn't think that all the $\pi$-related math on \emph{The Simpsons} 
is bogus. For example, in a ``Treehouse of Horror'' episode from October 30,
1995, we see the following equation, which, as we will learn shortly,
is true.

\begin{equation}\label{horrorequation}
e^{i\pi}=-1\\
\end{equation}

\begin{thm}\label{eulerformula}
For any $\theta\in\mathbb{R}$, we have $e^{i\theta}=\cos\theta+i\sin\theta$.
\end{thm}

\begin{proof}
Oh for crying out loud, this is so darned obvious I won't waste my time, but
if you must see a proof, please read \cite{LA1979} or any other good text on
complex analysis.\\
\end{proof}

\begin{cor}
The equation $e^{i\pi}=-1$ is true.
\end{cor}

\begin{proof}
If we let $\theta=\pi$ in Theorem \ref{eulerformula}, we obtain the following.
$$e^{i\pi}=\cos\pi+i\sin\pi=-1+i\cdot 0 = -1$$\\
\end{proof}

\begin{rem}
We note that Equation (\ref{horrorequation}) would look so much more beautiful
if 1 were added to both sides to give the following.
\begin{center}
\boxed{e^{i\pi}+1=0}
\end{center}
This single equation captures what many consider to be the five most important
numbers in mathematics. Good times.
\end{rem}
